\documentclass[10pt,a4paper,landscape]{article}

% ═══════════════════════════════════════════════════════════════════
% SADE MASTER SCREEN v4.0 - CORRETTO E VERIFICATO
% 5 pagine A4 landscape, una per pannello
% Tabelle verificate contro SADE.tex / SADE-MonstruousManual.tex
% ═══════════════════════════════════════════════════════════════════

\usepackage[T1]{fontenc}
\usepackage[utf8]{inputenc}
\usepackage{helvet}
\renewcommand{\familydefault}{\sfdefault}
\usepackage{geometry}
\usepackage[table]{xcolor}
\usepackage{colortbl}
\usepackage{booktabs}
\usepackage{longtable}
\usepackage{array}
\usepackage{multirow}
\usepackage{multicol}

% Layout A4 landscape
\geometry{
	a4paper,
	landscape,
	margin=1cm,
	top=1cm,
	bottom=1cm
}

% Colori
\definecolor{headerblue}{RGB}{41,82,163}
\definecolor{lightgray}{RGB}{245,245,245}

\pagestyle{empty}

% Spaziatura tabelle compatta
\renewcommand{\arraystretch}{1.0}
\setlength{\tabcolsep}{3pt}
\setlength{\columnsep}{10pt}
\setlength{\columnseprule}{0.3pt}

% Header pagina
\newcommand{\pageheader}[1]{%
	\noindent\colorbox{headerblue}{%
		\parbox{\dimexpr\textwidth-2\fboxsep\relax}{%
			\centering\color{white}\Large\bfseries #1
		}%
	}%
	\vspace{4pt}
}

\newcommand{\sectiontitle}[1]{%
	\vspace{3pt}
	\noindent{\normalsize\bfseries #1}
	\vspace{1pt}
}

\begin{document}

	% ═══════════════════════════════════════════════════════════════════
	% PAGINA 1: THAC0 E CLASSE ARMATURA
	% ═══════════════════════════════════════════════════════════════════

	\pageheader{PANNELLO 1: THAC0 E COMBATTIMENTO}

	\begin{multicols}{3}

		\sectiontitle{THAC0 Completo - Tutte le Classi}

		\scriptsize
		\noindent
		\begin{tabular}{c|cccccccccc}
			\toprule
			\rowcolor{lightgray}
			\textbf{Liv} & \textbf{1} & \textbf{2} & \textbf{3} & \textbf{4} & \textbf{5} & \textbf{6} & \textbf{7} & \textbf{8} & \textbf{9} & \textbf{10} \\
			\midrule
			Guerriero & 20 & 19 & 18 & 17 & 16 & 15 & 14 & 13 & 12 & 11 \\
			\rowcolor{lightgray}
			Sacerdote & 20 & 20 & 20 & 18 & 18 & 18 & 16 & 16 & 16 & 14 \\
			Canaglia & 20 & 19 & 19 & 18 & 17 & 17 & 16 & 15 & 15 & 14 \\
			\rowcolor{lightgray}
			Mago & 20 & 20 & 20 & 19 & 19 & 19 & 18 & 18 & 18 & 17 \\
			\bottomrule
		\end{tabular}

		\vspace{2pt}

		\noindent
		\begin{tabular}{c|cccccccccc}
			\toprule
			\rowcolor{lightgray}
			\textbf{Liv} & \textbf{11} & \textbf{12} & \textbf{13} & \textbf{14} & \textbf{15} & \textbf{16} & \textbf{17} & \textbf{18} & \textbf{19} & \textbf{20} \\
			\midrule
			Guerriero & 10 & 9 & 8 & 7 & 6 & 5 & 4 & 3 & 2 & 1 \\
			\rowcolor{lightgray}
			Sacerdote & 14 & 14 & 12 & 12 & 12 & 10 & 10 & 10 & 8 & 8 \\
			Canaglia & 13 & 13 & 12 & 12 & 11 & 11 & 10 & 10 & 9 & 9 \\
			\rowcolor{lightgray}
			Mago & 17 & 17 & 16 & 16 & 16 & 15 & 15 & 15 & 14 & 14 \\
			\bottomrule
		\end{tabular}

		\vspace{4pt}

		\sectiontitle{THAC0 delle Creature (per Dadi Vita)}

		\scriptsize
		\noindent
		\begin{tabular}{cc|cc}
			\toprule
			\rowcolor{lightgray}
			\textbf{DV} & \textbf{THAC0} & \textbf{DV} & \textbf{THAC0} \\
			\midrule
			$\leq$1/2 & 20 & 1 & 20 \\
			\rowcolor{lightgray}
			1+ & 19 & 2+ & 19 \\
			3+ & 17 & 4+ & 17 \\
			\rowcolor{lightgray}
			5+ & 15 & 6+ & 15 \\
			7+ & 13 & 8+ & 13 \\
			\rowcolor{lightgray}
			9+ & 11 & 10+ & 11 \\
			11+ & 9 & 12+ & 9 \\
			\rowcolor{lightgray}
			13+ & 7 & 14+ & 7 \\
			15+ & 5 & 16+ & 5 \\
			\bottomrule
		\end{tabular}

		\vspace{4pt}

		\sectiontitle{Formula Attacco}

		\footnotesize
		\begin{tabular}{l}
			\toprule
			\textbf{Tiro per Colpire:} 1d20 + mod. $\geq$ THAC0 $-$ CA bersaglio \\
			\midrule
			20 naturale = colpisce sempre \\
			1 naturale = manca sempre \\
			\bottomrule
		\end{tabular}

		\vspace{4pt}

		\sectiontitle{Tiro Critico}

		\footnotesize
		\begin{tabular}{l}
			\toprule
			Naturale 20, OPPURE \\
			naturale $\geq$18 con margine $\geq$5 \\
			\midrule
			Effetto: raddoppia i dadi danno arma \\
			\bottomrule
		\end{tabular}

		\columnbreak

		\sectiontitle{Classe Armatura (CA) Completa}

		\scriptsize
		\noindent
		\begin{tabular}{lc}
			\toprule
			\textbf{Tipo di Armatura} & \textbf{CA} \\
			\midrule
			Nessuna armatura & 10 \\
			\rowcolor{lightgray}
			Solo scudo & 9 \\
			Cuoio/imbottita & 8 \\
			\rowcolor{lightgray}
			Cuoio/imbottita + scudo & 7 \\
			Cuoio borch./anelli & 7 \\
			\rowcolor{lightgray}
			Cuoio borch./anelli + scudo & 6 \\
			Scaglie/pelle & 6 \\
			\rowcolor{lightgray}
			Scaglie/pelle + scudo & 5 \\
			Cotta di maglia & 5 \\
			\rowcolor{lightgray}
			Cotta di maglia + scudo & 4 \\
			Stecche/bande/bronzo & 4 \\
			\rowcolor{lightgray}
			Stecche/bande/bronzo + scudo & 3 \\
			Piastre & 3 \\
			\rowcolor{lightgray}
			Piastre + scudo & 2 \\
			Campale & 2 \\
			\rowcolor{lightgray}
			Campale + scudo & 1 \\
			Completa & 1 \\
			\rowcolor{lightgray}
			Completa + scudo & 0 \\
			\bottomrule
		\end{tabular}

		\vspace{4pt}

		\sectiontitle{Modificatori CA da Destrezza}

		\scriptsize
		\noindent
		\begin{tabular}{cc}
			\toprule
			\textbf{Destrezza} & \textbf{Mod. CA} \\
			\midrule
			1 & +5 \\
			\rowcolor{lightgray}
			2 & +5 \\
			3 & +4 \\
			\rowcolor{lightgray}
			4 & +3 \\
			5 & +2 \\
			\rowcolor{lightgray}
			6 & +1 \\
			7-14 & 0 \\
			\rowcolor{lightgray}
			15 & -1 \\
			16 & -2 \\
			\rowcolor{lightgray}
			17 & -3 \\
			18 & -4 \\
			\rowcolor{lightgray}
			19-20 & -4 \\
			21-22 & -5 \\
			\rowcolor{lightgray}
			23-24 & -6 \\
			25 & -7 \\
			\bottomrule
		\end{tabular}

		\vspace{4pt}

		\sectiontitle{Copertura e Occultamento}

		\scriptsize
		\noindent
		\begin{tabular}{lcc}
			\toprule
			\textbf{Bersaglio} & \textbf{Cop.} & \textbf{Occult.} \\
			\midrule
			Nascosto 25\% & -2 & -1 \\
			\rowcolor{lightgray}
			Nascosto 50\% & -4 & -2 \\
			Nascosto 75\% & -7 & -3 \\
			\rowcolor{lightgray}
			Nascosto 90\% & -10 & -4 \\
			\bottomrule
		\end{tabular}

		\columnbreak

		\sectiontitle{Modificatori Comuni al Colpire}

		\scriptsize
		\noindent
		\begin{tabular}{lc}
			\toprule
			\textbf{Situazione} & \textbf{Mod.} \\
			\midrule
			Attacco alle spalle & +2 \\
			\rowcolor{lightgray}
			Posizione sopraelevata (attaccante) & +1 \\
			Attaccante invisibile & +4 \\
			\rowcolor{lightgray}
			Difensore sorpreso & +1 \\
			Difensore prono/stordito & +4 \\
			\rowcolor{lightgray}
			Difensore indifeso (dorme, bloccato) & Colpito! \\
			Difensore invisibile & -4 \\
			\rowcolor{lightgray}
			Tiro a distanza, gittata media & -2 \\
			Tiro a distanza, gittata lunga & -5 \\
			\bottomrule
		\end{tabular}

		\vspace{4pt}

		\sectiontitle{Sorpresa}

		\footnotesize
		1d10 per parte coinvolta (segreto). \textbf{Soglia sorpresa: 1--3.} Chi è sorpreso perde il round; chi sorprende ottiene un round bonus d'azione libera.

		\scriptsize
		\noindent
		\begin{tabular}{lc}
			\toprule
			\textbf{Modificatore} & \textbf{al d10} \\
			\midrule
			Al buio / dorme & -4 \\
			\rowcolor{lightgray}
			Attaccato alle spalle & -2 \\
			In fuga / marcia forzata & -2 \\
			\rowcolor{lightgray}
			Invisibile o silenziato & -2 \\
			In scarsa luce & -1 \\
			\rowcolor{lightgray}
			Sospettoso / con cautela & +1/+2 \\
			Talento Allerta / è un Ranger & +1 \\
			\bottomrule
		\end{tabular}

		\vspace{4pt}

		\sectiontitle{Iniziativa (individuale)}

		\footnotesize
		Ogni PG/PNG/mostro tira \textbf{1d10} + mod. azione + velocità arma/tempo lancio $-$ mod. Destrezza. \textbf{Il totale PIÙ BASSO agisce per primo.}

		\scriptsize
		\noindent
		\begin{tabular}{lc}
			\toprule
			\textbf{Azione} & \textbf{Mod.} \\
			\midrule
			Muoversi / correre & +0 \\
			\rowcolor{lightgray}
			Alzarsi da terra & +5 \\
			Attacco mischia/distanza & +Vel.Arma \\
			\rowcolor{lightgray}
			Attacco senz'armi & +3 \\
			Caricare & +2 \\
			\rowcolor{lightgray}
			Lanciare incantesimo & +Tempo Lancio \\
			Bere pozione & +4 \\
			\rowcolor{lightgray}
			Sguainare arma & +2 \\
			Aprire porta & +3 \\
			\rowcolor{lightgray}
			Scacciare non-morti & +6 \\
			\bottomrule
		\end{tabular}

	\end{multicols}

	\clearpage

	% ═══════════════════════════════════════════════════════════════════
	% PAGINA 2: TIRI SALVEZZA
	% ═══════════════════════════════════════════════════════════════════

	\pageheader{PANNELLO 2: TIRI SALVEZZA E NON-MORTI}

	\sectiontitle{Scacciare Non-Morti (Sacerdote) --- risultato necessario su 2d10}

	\scriptsize
	\noindent
	\begin{tabular}{lcccccccccccc}
		\toprule
		\rowcolor{lightgray}
		& \multicolumn{12}{c}{\textbf{Livello del Sacerdote}} \\
		\rowcolor{lightgray}
		\textbf{Non-Morto} & \textbf{1} & \textbf{2} & \textbf{3} & \textbf{4} & \textbf{5} & \textbf{6} & \textbf{7} & \textbf{8} & \textbf{9} & \textbf{10-11} & \textbf{12-13} & \textbf{14+} \\
		\midrule
		Scheletro/1 DV & 10 & 7 & 4 & T & T & D & D & D* & D* & D* & D* & D* \\
		\rowcolor{lightgray}
		Zombie & 13 & 10 & 7 & 4 & T & T & D & D & D* & D* & D* & D* \\
		Ghoul/2 DV & 16 & 13 & 10 & 7 & 4 & T & T & D & D & D* & D* & D* \\
		\rowcolor{lightgray}
		Ombra/3-4 DV & 19 & 16 & 13 & 10 & 7 & 4 & T & T & D & D & D* & D* \\
		Wight/5 DV & 20 & 19 & 16 & 13 & 10 & 7 & 4 & T & T & D & D & D \\
		\rowcolor{lightgray}
		Ghast & -- & 20 & 19 & 16 & 13 & 10 & 7 & 4 & T & T & D & D \\
		Wraith/6 DV & -- & -- & 20 & 19 & 16 & 13 & 10 & 7 & 4 & T & T & D \\
		\rowcolor{lightgray}
		Mummia/7 DV & -- & -- & -- & 20 & 19 & 16 & 13 & 10 & 7 & 4 & T & T \\
		Spettro/8 DV & -- & -- & -- & -- & 20 & 19 & 16 & 13 & 10 & 7 & 4 & T \\
		\rowcolor{lightgray}
		Vampiro/9 DV & -- & -- & -- & -- & -- & 20 & 19 & 16 & 13 & 10 & 7 & 4 \\
		Fantasma/10 DV & -- & -- & -- & -- & -- & -- & 20 & 19 & 16 & 13 & 10 & 7 \\
		\rowcolor{lightgray}
		Lich/11+ DV & -- & -- & -- & -- & -- & -- & -- & 20 & 19 & 16 & 13 & 10 \\
		Speciali & -- & -- & -- & -- & -- & -- & -- & -- & 20 & 19 & 16 & 13 \\
		\bottomrule
	\end{tabular}
	\par
	{\scriptsize T = scaccia automaticamente; D = distrugge; D* = distrugge 2d4 aggiuntive; -- = fallisce automaticamente. Paladini scacciano come Sacerdoti di 2 livelli inferiori. Raggio 27m.}

	\vspace{4pt}

	\begin{multicols}{3}

		\sectiontitle{Tiri Salvezza - Guerriero}

		\scriptsize
		\noindent
		\begin{tabular}{c|ccccc}
			\toprule
			\rowcolor{lightgray}
			\textbf{Liv} & \textbf{P/V/M} & \textbf{Verga} & \textbf{Piet} & \textbf{Soffio} & \textbf{Inc} \\
			\midrule
			1-2 & 14 & 16 & 15 & 17 & 17 \\
			\rowcolor{lightgray}
			3-4 & 13 & 15 & 14 & 16 & 16 \\
			5-6 & 11 & 13 & 12 & 13 & 14 \\
			\rowcolor{lightgray}
			7-8 & 10 & 12 & 11 & 12 & 13 \\
			9-10 & 8 & 10 & 9 & 9 & 11 \\
			\rowcolor{lightgray}
			11-12 & 7 & 9 & 8 & 8 & 10 \\
			13-14 & 5 & 7 & 6 & 5 & 8 \\
			\rowcolor{lightgray}
			15-16 & 4 & 6 & 5 & 4 & 7 \\
			17-20 & 3 & 5 & 4 & 4 & 6 \\
			\bottomrule
		\end{tabular}

		\vspace{4pt}

		\sectiontitle{Tiri Salvezza - Sacerdote}

		\scriptsize
		\noindent
		\begin{tabular}{c|ccccc}
			\toprule
			\rowcolor{lightgray}
			\textbf{Liv} & \textbf{P/V/M} & \textbf{Verga} & \textbf{Piet} & \textbf{Soffio} & \textbf{Inc} \\
			\midrule
			1-3 & 10 & 14 & 13 & 16 & 15 \\
			\rowcolor{lightgray}
			4-6 & 9 & 13 & 12 & 15 & 14 \\
			7-9 & 7 & 11 & 10 & 13 & 12 \\
			\rowcolor{lightgray}
			10-12 & 6 & 10 & 9 & 12 & 11 \\
			13-15 & 5 & 9 & 8 & 11 & 10 \\
			\rowcolor{lightgray}
			16-18 & 4 & 8 & 7 & 10 & 9 \\
			19-20 & 2 & 6 & 5 & 8 & 7 \\
			\bottomrule
		\end{tabular}

		\columnbreak

		\sectiontitle{Tiri Salvezza - Canaglia}

		\scriptsize
		\noindent
		\begin{tabular}{c|ccccc}
			\toprule
			\rowcolor{lightgray}
			\textbf{Liv} & \textbf{P/V/M} & \textbf{Verga} & \textbf{Piet} & \textbf{Soffio} & \textbf{Inc} \\
			\midrule
			1-4 & 13 & 14 & 12 & 16 & 15 \\
			\rowcolor{lightgray}
			5-8 & 12 & 12 & 11 & 15 & 13 \\
			9-12 & 11 & 10 & 10 & 14 & 11 \\
			\rowcolor{lightgray}
			13-16 & 10 & 8 & 9 & 13 & 9 \\
			17-20 & 9 & 6 & 8 & 12 & 7 \\
			\bottomrule
		\end{tabular}

		\vspace{4pt}

		\sectiontitle{Tiri Salvezza - Mago}

		\scriptsize
		\noindent
		\begin{tabular}{c|ccccc}
			\toprule
			\rowcolor{lightgray}
			\textbf{Liv} & \textbf{P/V/M} & \textbf{Verga} & \textbf{Piet} & \textbf{Soffio} & \textbf{Inc} \\
			\midrule
			1-5 & 14 & 11 & 13 & 15 & 12 \\
			\rowcolor{lightgray}
			6-10 & 13 & 9 & 11 & 13 & 10 \\
			11-15 & 11 & 7 & 9 & 11 & 8 \\
			\rowcolor{lightgray}
			16-20 & 10 & 5 & 7 & 9 & 6 \\
			\bottomrule
		\end{tabular}

		\columnbreak

		\sectiontitle{Bonus Saggezza (vs magia, tutte le classi)}

		\scriptsize
		\noindent
		\begin{tabular}{lc}
			\toprule
			\textbf{Saggezza} & \textbf{Mod.} \\
			\midrule
			3 & -3 \\
			\rowcolor{lightgray}
			4 & -2 \\
			5-7 & -1 \\
			\rowcolor{lightgray}
			8-14 & 0 \\
			15 & +1 \\
			\rowcolor{lightgray}
			16 & +2 \\
			17 & +3 \\
			\rowcolor{lightgray}
			18-25 & +4 \\
			\bottomrule
		\end{tabular}

		\vspace{4pt}

		\sectiontitle{Bonus Razziali ai TS}

		\scriptsize
		\noindent
		\begin{tabular}{lc}
			\toprule
			\textbf{Nano/Halfling} & \\
			\midrule
			vs. Veleno & +5 \\
			\rowcolor{lightgray}
			vs. Bacchette & +3 \\
			vs. Incant./Magia & +3 \\
			\midrule
			\textbf{Gnomo} & \\
			\midrule
			vs. Bacchette & +3 \\
			\rowcolor{lightgray}
			vs. Incant./Magia & +3 \\
			\bottomrule
		\end{tabular}

	\end{multicols}

	\clearpage

	% ═══════════════════════════════════════════════════════════════════
	% PAGINA 3: DANNI ARMI E MODIFICATORI
	% ═══════════════════════════════════════════════════════════════════

	\pageheader{PANNELLO 3: DANNO ARMI E MODIFICATORI}

	\begin{multicols}{3}

		\sectiontitle{Danni Armi da Mischia}

		\scriptsize
		\noindent
		\begin{tabular}{lcccc}
			\toprule
			\rowcolor{lightgray}
			\textbf{Arma} & \textbf{Costo} & \textbf{Peso} & \textbf{vs M} & \textbf{vs G} \\
			\midrule
			Ascia da battaglia & 5 mo & 3,5 kg & 1d8 & 1d8 \\
			\rowcolor{lightgray}
			Ascia da lancio & 1 mo & 2 kg & 1d6 & 1d4 \\
			Ascia da mano & 1 mo & 2,5 kg & 1d6 & 1d4 \\
			\rowcolor{lightgray}
			Ascia bipenne & 10 mo & 7,5 kg & 1d8 & 1d8 \\
			Bastone ferrato & 1 mo & 2 kg & 1d6 & 1d6 \\
			\rowcolor{lightgray}
			Clava & --- & 1,5 kg & 1d6 & 1d3 \\
			Daga/Pugnale & 2 mo & 0,5 kg & 1d4 & 1d3 \\
			\rowcolor{lightgray}
			Falce & 6 mo & 5 kg & 1d4+1 & 1d4 \\
			Falcione & 8 mo & 5 kg & 2d4 & 1d6+1 \\
			\rowcolor{lightgray}
			Lancia & 8 ma & 2,5 kg & 1d6 & 1d8 \\
			Lancia da cavaliere & 15 mo & 5 kg & 1d8 & 3d6* \\
			\rowcolor{lightgray}
			Mazza & 8 mo & 4 kg & 1d6+1 & 1d6 \\
			Mazza grande & 5 mo & 5 kg & 1d6+1 & 1d6 \\
			\rowcolor{lightgray}
			Martello da guerra & 1 mo & 3 kg & 1d4+1 & 1d4 \\
			Mazzafrusto & 15 mo & 7,5 kg & 1d6+1 & 2d4 \\
			\rowcolor{lightgray}
			Morning star & 10 mo & 6 kg & 2d4 & 1d6+1 \\
			Piccone militare & 8 mo & 3 kg & 1d6+1 & 2d4 \\
			\rowcolor{lightgray}
			Picca & 5 mo & 6 kg & 1d6 & 1d8 \\
			Randello & --- & 1,5 kg & 1d6 & 1d3 \\
			\rowcolor{lightgray}
			Scimitarra & 15 mo & 2 kg & 1d8 & 1d8 \\
			Spada a due mani & 50 mo & 7,5 kg & 1d10 & 3d6 \\
			\rowcolor{lightgray}
			Spada bastarda & 25 mo & 5 kg & 2d4 & 2d8 \\
			Spada corta & 10 mo & 1,5 kg & 1d6 & 1d8 \\
			\rowcolor{lightgray}
			Spada lunga & 15 mo & 2 kg & 1d8 & 1d12 \\
			Tridente & 15 mo & 2,5 kg & 1d6+1 & 3d4 \\
			\bottomrule
		\end{tabular}
		\par
		\noindent\parbox{\linewidth}{\scriptsize *Solo in carica a cavallo. M=vs creature medie o pi\`u piccole; G=vs creature grandi.}

		\columnbreak

		\sectiontitle{Danni Armi da Distanza}

		\scriptsize
		\noindent
		\begin{tabular}{lcccc}
			\toprule
			\rowcolor{lightgray}
			\textbf{Arma} & \textbf{Costo} & \textbf{Peso} & \textbf{vs M} & \textbf{vs G} \\
			\midrule
			Arco corto & 30 mo & 1 kg & 1d6 & 1d6 \\
			\rowcolor{lightgray}
			Arco lungo & 75 mo & 1,5 kg & 1d6 & 1d6 \\
			Arco composito corto & 75 mo & 1 kg & 1d6 & 1d6 \\
			\rowcolor{lightgray}
			Arco composito lungo & 100 mo & 1,5 kg & 1d6 & 1d6 \\
			Balestra leggera & 35 mo & 3,5 kg & 1d4 & 1d4 \\
			\rowcolor{lightgray}
			Balestra pesante & 50 mo & 7 kg & 1d4+1 & 1d6+1 \\
			Fionda & 5 mr & --- & 1d4 & 1d4 \\
			\rowcolor{lightgray}
			Giavellotto & 5 ma & 1 kg & 1d6 & 1d6 \\
			\bottomrule
		\end{tabular}

		\vspace{6pt}

		\sectiontitle{Modificatori Danno da Forza}

		\scriptsize
		\noindent
		\begin{tabular}{lcc}
			\toprule
			\rowcolor{lightgray}
			\textbf{Forza} & \textbf{Attacco} & \textbf{Danno} \\
			\midrule
			3 & -3 & -1 \\
			\rowcolor{lightgray}
			4-5 & -2 & -1 \\
			6-7 & 0 & -1 \\
			\rowcolor{lightgray}
			8-15 & 0 & 0 \\
			16 & 0 & +1 \\
			\rowcolor{lightgray}
			17 & +1 & +1 \\
			18 & +1 & +2 \\
			\rowcolor{lightgray}
			18/01-50 & +1 & +3 \\
			18/51-75 & +2 & +3 \\
			\rowcolor{lightgray}
			18/76-90 & +2 & +4 \\
			18/91-99 & +2 & +5 \\
			\rowcolor{lightgray}
			18/00 & +3 & +6 \\
			19 & +3 & +7 \\
			\rowcolor{lightgray}
			20 & +3 & +8 \\
			21-22 & +4 & +9/+10 \\
			\rowcolor{lightgray}
			23-25 & +5-7 & +11-14 \\
			\bottomrule
		\end{tabular}

		\columnbreak

		\sectiontitle{Modificatori da Costituzione}

		\scriptsize
		\noindent
		\begin{tabular}{lc}
			\toprule
			\rowcolor{lightgray}
			\textbf{Costituzione} & \textbf{PF/Dado} \\
			\midrule
			3 & -2 \\
			\rowcolor{lightgray}
			4-6 & -1 \\
			7-14 & 0 \\
			\rowcolor{lightgray}
			15 & +1 \\
			16 & +2 \\
			\rowcolor{lightgray}
			17 & +2 (+3 Guerrieri) \\
			18 & +2 (+4 Guerrieri) \\
			\rowcolor{lightgray}
			19-20 & +2 (+5 Guerrieri) \\
			21-23 & +2 (+6 Guerrieri) \\
			\rowcolor{lightgray}
			24-25 & +2 (+7 Guerrieri) \\
			\bottomrule
		\end{tabular}

		\vspace{4pt}

		\sectiontitle{Morte e Guarigione}

		\scriptsize
		\noindent
		\begin{tabular}{ll}
			\toprule
			\rowcolor{lightgray}
			\textbf{PF} & \textbf{Condizione} \\
			\midrule
			10+ PF & Sano \\
			\rowcolor{lightgray}
			1-9 PF & Ferito \\
			0 PF & Incosciente, stabile \\
			\rowcolor{lightgray}
			-1 a -9 PF & Morente (-1 PF/round) \\
			-10 PF & Morto \\
			\bottomrule
		\end{tabular}

		\vspace{4pt}

		\noindent
		\begin{tabular}{ll}
			\toprule
			\rowcolor{lightgray}
			\textbf{Recupero} & \textbf{PF} \\
			\midrule
			Riposo (8 ore) & 1 PF/giorno \\
			\rowcolor{lightgray}
			Riposo completo + cure & 1d3 PF/giorno \\
			Attività faticosa & Nessuno \\
			\bottomrule
		\end{tabular}

	\end{multicols}

	\clearpage

	% ═══════════════════════════════════════════════════════════════════
	% PAGINA 4: INCANTESIMI AL GIORNO
	% ═══════════════════════════════════════════════════════════════════

	\pageheader{PANNELLO 4: INCANTESIMI AL GIORNO}

	\begin{multicols}{3}

		\sectiontitle{Mago}

		\scriptsize
		\noindent
		\begin{tabular}{c|ccccccccc}
			\toprule
			\rowcolor{lightgray}
			\textbf{L} & \textbf{1} & \textbf{2} & \textbf{3} & \textbf{4} & \textbf{5} & \textbf{6} & \textbf{7} & \textbf{8} & \textbf{9} \\
			\midrule
			1 & 1 & - & - & - & - & - & - & - & - \\
			\rowcolor{lightgray}
			2 & 2 & - & - & - & - & - & - & - & - \\
			3 & 2 & 1 & - & - & - & - & - & - & - \\
			\rowcolor{lightgray}
			4 & 3 & 2 & - & - & - & - & - & - & - \\
			5 & 4 & 2 & 1 & - & - & - & - & - & - \\
			\rowcolor{lightgray}
			6 & 4 & 2 & 2 & - & - & - & - & - & - \\
			7 & 4 & 3 & 2 & 1 & - & - & - & - & - \\
			\rowcolor{lightgray}
			8 & 4 & 3 & 3 & 2 & - & - & - & - & - \\
			9 & 4 & 3 & 3 & 2 & 1 & - & - & - & - \\
			\rowcolor{lightgray}
			10 & 4 & 4 & 3 & 2 & 2 & - & - & - & - \\
			11 & 4 & 4 & 4 & 3 & 3 & - & - & - & - \\
			\rowcolor{lightgray}
			12 & 4 & 4 & 4 & 4 & 4 & 1 & - & - & - \\
			13 & 5 & 5 & 5 & 4 & 4 & 2 & - & - & - \\
			\rowcolor{lightgray}
			14 & 5 & 5 & 5 & 4 & 4 & 2 & 1 & - & - \\
			15 & 5 & 5 & 5 & 5 & 5 & 2 & 1 & - & - \\
			\rowcolor{lightgray}
			16 & 5 & 5 & 5 & 5 & 5 & 3 & 2 & 1 & - \\
			17 & 5 & 5 & 5 & 5 & 5 & 3 & 3 & 2 & - \\
			\rowcolor{lightgray}
			18 & 5 & 5 & 5 & 5 & 5 & 3 & 3 & 2 & 1 \\
			19 & 5 & 5 & 5 & 5 & 5 & 3 & 3 & 3 & 1 \\
			\rowcolor{lightgray}
			20 & 5 & 5 & 5 & 5 & 5 & 4 & 3 & 3 & 2 \\
			\bottomrule
		\end{tabular}

		\columnbreak

		\sectiontitle{Chierico e Druido (progressione condivisa)}

		\scriptsize
		\noindent
		\begin{tabular}{c|ccccccc}
			\toprule
			\rowcolor{lightgray}
			\textbf{L} & \textbf{1} & \textbf{2} & \textbf{3} & \textbf{4} & \textbf{5} & \textbf{6*} & \textbf{7**} \\
			\midrule
			1 & 1 & - & - & - & - & - & - \\
			\rowcolor{lightgray}
			2 & 2 & - & - & - & - & - & - \\
			3 & 2 & 1 & - & - & - & - & - \\
			\rowcolor{lightgray}
			4 & 3 & 2 & - & - & - & - & - \\
			5 & 3 & 3 & 1 & - & - & - & - \\
			\rowcolor{lightgray}
			6 & 3 & 3 & 2 & - & - & - & - \\
			7 & 3 & 3 & 2 & 1 & - & - & - \\
			\rowcolor{lightgray}
			8 & 3 & 3 & 3 & 2 & - & - & - \\
			9 & 4 & 4 & 3 & 2 & 1 & - & - \\
			\rowcolor{lightgray}
			10 & 4 & 4 & 3 & 3 & 2 & - & - \\
			11 & 5 & 4 & 4 & 3 & 2 & 1 & - \\
			\rowcolor{lightgray}
			12 & 6 & 5 & 5 & 3 & 2 & 2 & - \\
			13 & 6 & 6 & 6 & 4 & 2 & 2 & - \\
			\rowcolor{lightgray}
			14 & 6 & 6 & 6 & 5 & 3 & 2 & 1 \\
			15 & 6 & 6 & 6 & 6 & 4 & 2 & 1 \\
			\rowcolor{lightgray}
			16 & 7 & 7 & 7 & 6 & 4 & 3 & 1 \\
			17 & 7 & 7 & 7 & 7 & 5 & 3 & 2 \\
			\rowcolor{lightgray}
			18 & 8 & 8 & 8 & 8 & 6 & 4 & 2 \\
			19 & 9 & 9 & 8 & 8 & 6 & 4 & 2 \\
			\rowcolor{lightgray}
			20 & 9 & 9 & 9 & 8 & 7 & 5 & 2 \\
			\bottomrule
		\end{tabular}
		\par
		\noindent\parbox{\linewidth}{\scriptsize *Sag 17+. **Sag 18+. Chierico e Druido condividono questa progressione in SADE (differiscono per titoli e PE).}

		\columnbreak

		\sectiontitle{Bardo (+ uso pergamene)}

		\scriptsize
		\noindent
		\begin{tabular}{c|cccccc|c}
			\toprule
			\rowcolor{lightgray}
			\textbf{L} & \textbf{1} & \textbf{2} & \textbf{3} & \textbf{4} & \textbf{5} & \textbf{6} & \textbf{Perg.} \\
			\midrule
			1 & 1 & - & - & - & - & - & 20\% \\
			\rowcolor{lightgray}
			2 & 2 & - & - & - & - & - & 20\% \\
			3 & 2 & 1 & - & - & - & - & 30\% \\
			\rowcolor{lightgray}
			4 & 3 & 1 & - & - & - & - & 30\% \\
			5 & 3 & 2 & - & - & - & - & 40\% \\
			\rowcolor{lightgray}
			6 & 3 & 2 & 1 & - & - & - & 40\% \\
			7 & 3 & 3 & 1 & - & - & - & 50\% \\
			\rowcolor{lightgray}
			8 & 3 & 3 & 2 & - & - & - & 60\% \\
			9 & 3 & 3 & 2 & 1 & - & - & 70\% \\
			\rowcolor{lightgray}
			10 & 3 & 3 & 3 & 1 & - & - & 80\% \\
			11 & 3 & 3 & 3 & 2 & - & - & 90\% \\
			\rowcolor{lightgray}
			12 & 3 & 3 & 3 & 2 & 1 & - & 91\% \\
			13 & 3 & 3 & 3 & 3 & 1 & - & 91\% \\
			\rowcolor{lightgray}
			14 & 3 & 3 & 3 & 3 & 2 & - & 92\% \\
			15 & 4 & 3 & 3 & 3 & 2 & 1 & 92\% \\
			\rowcolor{lightgray}
			16 & 4 & 4 & 3 & 3 & 3 & 1 & 93\% \\
			17 & 4 & 4 & 4 & 3 & 3 & 2 & 93\% \\
			\rowcolor{lightgray}
			18 & 4 & 4 & 4 & 4 & 3 & 2 & 94\% \\
			19 & 4 & 4 & 4 & 4 & 4 & 3 & 94\% \\
			\rowcolor{lightgray}
			20 & 4 & 4 & 4 & 4 & 4 & 3 & 95\% \\
			\bottomrule
		\end{tabular}

		\vspace{6pt}

		\sectiontitle{Regole Lancio Incantesimi}

		\scriptsize
		\noindent
		\begin{tabular}{ll}
			\toprule
			\rowcolor{lightgray}
			\textbf{Situazione} & \textbf{Effetto} \\
			\midrule
			Dichiarazione & In iniziativa, prima di altri \\
			\rowcolor{lightgray}
			Danno subito & TS Incantesimi o perso \\
			Movimento & Interrompe il lancio \\
			\rowcolor{lightgray}
			Silenziato/Legato & Niente verbale/somatico \\
			Senza materiali & Impossibile da lanciare \\
			\bottomrule
		\end{tabular}

	\end{multicols}

	\clearpage

	% ═══════════════════════════════════════════════════════════════════
	% PAGINA 5: ESPLORAZIONE E INCONTRI
	% ═══════════════════════════════════════════════════════════════════

	\pageheader{PANNELLO 5: ESPLORAZIONE E INCONTRI}

	\begin{multicols}{3}

		\sectiontitle{Ingombro e Movimento (in metri)}

		\scriptsize
		\noindent
		\begin{tabular}{lccc}
			\toprule
			\rowcolor{lightgray}
			\textbf{Ingombro} & \textbf{Base 12} & \textbf{Base 9} & \textbf{Base 6} \\
			\midrule
			Leggero & 12 & 9 & 6 \\
			\rowcolor{lightgray}
			Moderato & 8 & 6 & 4 \\
			Pesante & 6 & 4 & 3 \\
			\rowcolor{lightgray}
			Grave & 4 & 3 & 2 \\
			\bottomrule
		\end{tabular}
		\par
		\noindent\parbox{\linewidth}{\scriptsize Base 12: Umano/Elfo/Mezzelfo; Base 9: Nano/Gnomo/Semiorco; Base 6: Halfling. Categoria in base al peso trasportato/Forza (Tab. Ingombro).}

		\vspace{4pt}

		\sectiontitle{Velocità Esplorazione Dungeon}

		\scriptsize
		\noindent
		\begin{tabular}{lcc}
			\toprule
			\rowcolor{lightgray}
			\textbf{Tipo} & \textbf{m/turno} & \textbf{Note} \\
			\midrule
			Cauta & 18 & Mappa automatica \\
			\rowcolor{lightgray}
			Standard & 36 & Mappa possibile \\
			Veloce & 72 & Nessuna ricerca \\
			\bottomrule
		\end{tabular}

		\vspace{4pt}

		\sectiontitle{Porte Segrete (1d6) / Sentire Rumori (1d6)}

		\scriptsize
		\noindent
		\begin{tabular}{lc}
			\toprule
			\rowcolor{lightgray}
			\textbf{Situazione} & \textbf{Succ.} \\
			\midrule
			Umano cerca attivamente & 1-2 \\
			\rowcolor{lightgray}
			Elfo passa vicino / cerca & 1-2 / 1-4 \\
			Nano (costruzioni nane) & 1-3 \\
			\rowcolor{lightgray}
			Sentire rumori (Umano) & 1 \\
			Sentire rumori (Elfo/Nano/Half.) & 1-2 \\
			\bottomrule
		\end{tabular}

		\columnbreak

		\sectiontitle{Abilità Ladro Base (Livello 1)}

		\scriptsize
		\noindent
		\begin{tabular}{lc}
			\toprule
			\rowcolor{lightgray}
			\textbf{Abilità} & \textbf{\%} \\
			\midrule
			Scassinare Serrature & 30\% \\
			\rowcolor{lightgray}
			Scovare/Rimuovere Trappole & 25\% \\
			Borseggiare & 30\% \\
			\rowcolor{lightgray}
			Muoversi Silenziosamente & 20\% \\
			Nascondersi nelle Ombre & 15\% \\
			\rowcolor{lightgray}
			Sentire Rumori & 15\% \\
			Scalare Pareti & 60\% \\
			\rowcolor{lightgray}
			Leggere Linguaggi & 0\% (dal liv.4) \\
			\bottomrule
		\end{tabular}

		\vspace{4pt}

		\sectiontitle{Pugnalata alle Spalle}

		\scriptsize
		\noindent
		\begin{tabular}{lcc}
			\toprule
			\rowcolor{lightgray}
			\textbf{Livello Ladro} & \textbf{Bonus} & \textbf{Molt.} \\
			\midrule
			1-4 & +4 & x2 \\
			\rowcolor{lightgray}
			5-8 & +4 & x3 \\
			9-12 & +4 & x4 \\
			\rowcolor{lightgray}
			13+ & +4 & x5 \\
			\bottomrule
		\end{tabular}

		\vspace{4pt}

		\sectiontitle{Incontri Casuali / Distanza}

		\scriptsize
		\noindent
		\begin{tabular}{lcc}
			\toprule
			\rowcolor{lightgray}
			\textbf{Ambiente} & \textbf{Frequenza} & \textbf{Distanza} \\
			\midrule
			Dungeon & 1d20/3 turni & 2d6 x 3m \\
			\rowcolor{lightgray}
			Wilderness aperto & 1d20/giorno & 4d6 x 9m \\
			Foresta/Giungla & 1d20/giorno & 2d6 x 3m \\
			\rowcolor{lightgray}
			Città & 1d20/6 turni & 2d6 x 3m \\
			\bottomrule
		\end{tabular}
		\par
		\noindent\parbox{\linewidth}{\scriptsize Successo su 1-2 (su 1d20) salvo indicazioni diverse del GM.}

		\columnbreak

		\sectiontitle{Sorpresa (1d10 per parte)}

		\scriptsize
		\noindent
		\begin{tabular}{lc}
			\toprule
			\rowcolor{lightgray}
			\textbf{Risultato} & \textbf{Effetto} \\
			\midrule
			1-3 & Sorpreso (perde il round) \\
			\rowcolor{lightgray}
			4+ & Non sorpreso \\
			\bottomrule
		\end{tabular}
		\par
		\noindent\parbox{\linewidth}{\scriptsize Modificatori situazionali: v. Pannello 1.}

		\vspace{4pt}

		\sectiontitle{Tabella Reazioni (2d6)}

		\scriptsize
		\noindent
		\begin{tabular}{cl}
			\toprule
			\rowcolor{lightgray}
			\textbf{2d6} & \textbf{Reazione} \\
			\midrule
			2 & Ostile, attacca subito \\
			\rowcolor{lightgray}
			3-5 & Ostile, potrebbe attaccare \\
			6-8 & Cauto, incerto \\
			\rowcolor{lightgray}
			9-11 & Neutrale, negoziabile \\
			12 & Amichevole \\
			\bottomrule
		\end{tabular}
		\par
		\noindent\parbox{\linewidth}{\scriptsize Mod.: Carisma da -3 (Car 3) a +3 (Car 18).}

		\vspace{4pt}

		\sectiontitle{Grado di Morale (2d10 vs punteggio)}

		\scriptsize
		\noindent
		\begin{tabular}{lc}
			\toprule
			\rowcolor{lightgray}
			\textbf{Grado} & \textbf{Punt.} \\
			\midrule
			Inaffidabile & 2-4 \\
			\rowcolor{lightgray}
			Instabile & 5-7 \\
			Medio & 8-10 \\
			\rowcolor{lightgray}
			Fermo & 11-12 \\
			Élite & 13-14 \\
			\rowcolor{lightgray}
			Campione & 15-16 \\
			Fanatico & 17-18 \\
			\rowcolor{lightgray}
			Senza Paura & 19-20 \\
			\bottomrule
		\end{tabular}
		\par
		\noindent\parbox{\linewidth}{\scriptsize Tira 2d10 $\leq$ punteggio = resiste. Prova quando: primo danno subito, 50\% gruppo caduto, leader sconfitto, potere soverchiante.}

	\end{multicols}

	\vfill

	\begin{center}
		\small
		\textcolor{headerblue}{\textbf{SADE Master Screen v4.0}} • Sistema Avanzato Dungeons \& Dragons Espansioni • \copyright{} 2024 Andres Zanzani
	\end{center}

\end{document}
